% =============================================================================
% Section 6: Discussion
% =============================================================================

\subsection{Capability-First Audit: Ontological Implications}\label{subsec:ontological-implications}

The empirical results support the event-first design; the deeper contribution lies in four ontological questions.

\textbf{Q1: What is the ontological status of a capability?} In ADL Lite, a capability is a generically dependent continuant that exists only from its genesis event (\texttt{REGISTER}) onward; prior to that, there is only potential. This parallels BFO's view that a process exists only from its beginning.

\textbf{Q2: Are status and confidence real or derived?} Status is computed by $\delta(C)$ and has no existence independent of the chain, aligning with DOLCE's view that some properties are constructed by cognitive agents rather than discovered~\cite{dolce}.

\textbf{Q3: What happens to a capability when it is deprecated?} Deprecation appends a \texttt{DEPRECATE} event; the capability retains its \texttt{concept\_id} but its status becomes \texttt{deprecated}, preserving a complete audit trail. The entity is not destroyed---its identity persists---but its audit status changes, reflecting the distinction between identity and status central to the event-first design.

\textbf{Q4: What is the ontological cost of the event-first stance?} Historical queries are natural, but snapshot queries require scanning the chain: computing all currently \texttt{validated} concepts requires evaluating $\delta(C)$ per chain, which is $O(n)$. This is deliberate---event-first trades query-time computation for the elimination of stored mutable state and its consistency defects---and the complexity is linear, predictable, and acceptable at the target scale ($10^5$--$10^6$ events).

\textbf{Space-time trade-off.} The append-only design trades storage growth for query-time recomputation. At $5\times10^5$ events, zstd+msgpack cold storage achieves 8.4$\times$ compression over JSONL (E27), projected to exceed 10$\times$ at $10^6$ events. The warm index caches derived snapshots for $O(1)$ incremental queries, so the $O(n)$ cost is incurred only on cold starts or cache invalidations; secondary indexing could be introduced without changing canonical derivation semantics.

\subsection{Capability Registry Contribution}\label{subsec:conceptual-modeling}

The three contributions (Section~\ref{subsec:contributions})---ontological analysis, formal semantics, and architectural verification---position ADL Lite as a lightweight audit layer.

\textbf{Ontological analysis.} The cross-mapping to BFO, DOLCE, and UFO (Table~\ref{tab:upper-ontology-mapping}) provides methodological guidance for event-first systems. The two-level account (\S\ref{subsec:two-level-account}) resolves the process-vs.-record ambiguity of event-sourced systems.

\textbf{Formal semantics.} Eight machine-verified properties in Coq 8.18.0 (Theorems~1--7, 9) and the Python-level property (Theorem~8) provide formal assurance of deterministic, auditable state derivation. The decidable precondition language (Theorem~8) provides the first $O(1)$-per-rule evaluation guarantee for capability-lifecycle audit.

\textbf{Architectural verification.} Validation on 201~EventChains (9,300~events) confirms linear scalability and zero integrity failures, whereas the head-to-head benchmark (E19, Table~\ref{tab:e19-benchmark}) demonstrates competitive audit cost against nanopublications, PROV-O, and Git-only baselines.

The complete OWL~2~DL module and SHACL shape graph provide machine-checkable structural validation for the applied ontology community, bridging the gap between event-sourced runtime semantics and declarative ontology tooling.

\subsection{Limitations}\label{subsec:limitations}

The results must be interpreted in light of the following limitations.

\textbf{L1. Informal foundational alignment.} The ontological analysis is conceptual, relying on expert judgment; it is not validated with automated alignment tools (e.g., LogMap, AML). The event-first stance is a design commitment, not an ontological necessity; alternative BFO/IAO-compatible approaches are theoretically valid.~(FW1)

\textbf{L2. Append-only accumulation.} The append-only design leads to junk accumulation; feasibility at $10^5$--$10^6$ events was confirmed empirically (E19: 133.9~s at 14,938 events/s).~(FW2)

\textbf{L3. Collusion vulnerability.} A single actor can inflate $\gamma(C)$ to $0.99$ through repeated self-validation (E14); calibration lacks game-theoretic guarantees. $N_{\min} \geq 2$ prevents a single actor from triggering \textit{validated}, but does not address Sybil attacks. Phase~1.5 (implemented, pre-release) adds \texttt{did:key} Ed25519 signatures; full Sybil resistance requires Phase~3 staking.~(FW4, FW9)

\textbf{L4. Cryptographic authentication gap.} The implementation relied on self-declared string identifiers; a minimal \texttt{did:key} layer was implemented (\S\ref{subsec:trust-model}) but is a Phase~1.5 pre-implementation, not production-grade.~(FW4)

\textbf{L5. Absence of expert ground truth.} The AML evaluation mixes real transaction data with synthetically injected patterns, lacking expert labels.~(FW5)

\textbf{L6. Unstructured agent communication.} L2 Markdown imposes no structural constraints on agent reasoning; template-governed sections were evaluated in E20 with modest gains, and full L2 narrative validation remains future work.~(FW6)

\textbf{L7. Brittle capability discovery.} \texttt{DataImporter.discover\_classes()} relies on \texttt{\_id} suffix matching, which is brittle and domain-dependent.~(FW7)

\textbf{L8. Semantic Web interoperability.} OWL~2 DL bidirectional import/export, RDF-star export, PROV-O mapping, JSON-LD serialisation, and SHACL validation were implemented as of v0.6.0-alpha (Appendix~\ref{app:implementation}); SHACL provides machine-checkable validation via six shapes.

\textbf{L9. No incentive-compatible validation.} $\gamma(C)$ aggregates self-reported confidence without staking or slashing; linear calibration weights confidence by per-actor accuracy but lacks game-theoretic guarantees.~(FW9)

\textbf{L10. Incomplete machine-checked proofs.} Eight theorems (T1--T7, T9) were machine-verified in Coq 8.18.0 with zero remaining \texttt{Admitted} lemmas; T8 is a Python-level property. TLA$^{+}$ was model-checked for bounded chains. Three previously stubbed axioms were replaced with full definitions; six axioms remain simplified to \texttt{True}---these stubs do not weaken the algebraic proofs but do not rule out adversarial scenarios.~(FW12)

\textbf{L11. Overstatement risk in audit claims.} Readers may overgeneralize tamper-evidence guarantees to adversarial settings; we explicitly stated that the SHA-256 chain detects tampering post-hoc but does not prevent it.~(FW4)

\textbf{L12. Fork resolution.} Fork resolution uses CRDT lattice semantics (LUB); for multi-way concurrent merges with conflicting lifecycle events, semantic merge policies are planned for Phase~3.~(FW4)

\textbf{L13. Multi-way CRDT merge evaluation.} E26/E31/E28 validated merge correctness, sub-millisecond merge latency, and 10{,}000-agent integrity; v0.6.0-alpha introduced N$\geq$3 chain merging via \texttt{functools.reduce}. Large-scale distributed evaluation with WAN latency remains future work.~(FW2)

\textbf{L14. No human expert validation.} The AML case study and multi-agent simulations used synthetic validators and author-internal scoring; E32 is a simulated proxy, not a substitute for independent human review (E5, FW15).~(FW5)

\textbf{L15. No formal Merkle log comparison at scale.} E33 compared ADL Lite's linear chain with Merkle baselines on a single node; distributed Merkle-log benchmarks are future work.~(FW2)

\textbf{L16. OWL/SHACL completeness.} $\delta$ and the precondition language were intentionally excluded from the OWL fragment because they exceed DL expressivity; the formal semantics are captured in the Coq development, not the OWL fragment.~(FW1)

\subsection{Security Recommendations for Early Adopters}\label{subsec:security-recommendations}

The adversarial trials (E4e, Section~\ref{subsec:adversarial-trials}) and the Phase~1--3 authentication roadmap (Appendix~K) motivate four recommendations, all implementable with the current v0.6.0-alpha codebase. (1)~\textbf{Key signing per actor:} each actor should use a dedicated Ed25519 key pair in the \texttt{KeyRegistry} (\texttt{sign\_event()}/\texttt{verify\_signature()} wrappers attach W3C Linked Data Proofs to each event; unsigned events are rejected by the ActionExecutor), mitigating impersonation (B3, Table~\ref{tab:undetectable-attacks})---the single most effective Phase~1.5 hardening measure. (2)~\textbf{Minimum-validator policies:} \texttt{min\_distinct\_validators()} (default 1, configurable via \texttt{adl\_core\_ontology.yaml}) should be set to $N_{\min} \geq 2$ for \textsc{Validate} and $N_{\min} \geq 3$ for \textsc{Deprecate}, preventing the single-actor collusion of E14 at append time. (3)~\textbf{Repository tamper detection:} store EventChains in Git with \texttt{reflog} enabled so forced pushes and rebases (B2) remain detectable; for stronger guarantees, batch-anchor Merkle roots to an external transparency log (e.g., Sigstore Rekor~\cite{sigstore2022}) via \texttt{TransparencyAnchor}---the E33 comparison (Section~\ref{subsec:e33}) shows Merkle inclusion proofs add $<$1~ms verification latency at $O(\log n)$ proof size. (4)~\textbf{Scope enforcement via ACL prefixes:} use the scope prefix system (\texttt{public/}, \texttt{private/<org>}, \texttt{user/<id>}, \texttt{shared/<collab>}) with \texttt{ADLValidator} scope-ACL rules and Git branch protection to prevent unauthorized writes. Together these reduce the attack surface from the open-collaboration model to a threshold-guarded model where compromise requires multiple colluding actors or repository-level access; they do not eliminate the need for Phase~3 economic staking (FW9).

\subsection{Formal Adversarial Robustness Framework}\label{subsec:adversarial-framework}

The adversarial trials (E4e) and baseline comparison (E37) can be unified under a formal robustness framework that makes the trust-model boundary explicit. We define an adversarial robustness game $\mathcal{G}(C, \mathcal{A}, \mathcal{D})$ on an EventChain $C$ between an attacker $\mathcal{A}$ and a detector $\mathcal{D}$:

\begin{enumerate}[label=(\arabic*)]
    \item $\mathcal{A}$ selects an attack class $\alpha$ from the taxonomy in Appendix~C and applies it to $C$, producing a modified chain $C'$.
    \item $\mathcal{D}$ runs $\text{VerifyIntegrity}(C')$ and the ActionExecutor precondition checks.
    \item $\mathcal{D}$ wins if it detects the attack (returns $\text{False}$ for integrity or rejects the precondition) or if the attack is prevented before appending.
\end{enumerate}

The framework yields three \emph{robustness zones}. (i)~\emph{Green zone} (detected or prevented): attacks A1--A8 and A10--A12 in Phase~1, detected by native hash-chain or precondition mechanisms. (ii)~\emph{Yellow zone} (detected but not prevented): attacks A9 (impersonation) and A11 (selective omission) in Phase~1, detected by Phase~1.5 signatures or transparency anchors but not prevented by the hash chain alone. (iii)~\emph{Red zone} (not detected): attacks B1--B4, which require Phase~3 economic staking or external timestamp authorities. This three-zone taxonomy maps directly to the baseline comparison (Supplementary Table S4): ADL Lite's green zone is larger than nanopublications' (structural ordering via hash chaining) and comparable to PROV-O's and Git-only's, but its yellow zone is smaller because lifecycle preconditions are native rather than external.

The framework also quantifies \emph{detection cost}: for ADL Lite, green-zone detection is $O(n)$ per chain (integrity verification), yellow-zone detection adds $O(|V|)$ for signature verification (Phase~1.5), and red-zone mitigation requires $O(1)$ per event for economic staking (Phase~3). The baseline comparison reveals that nanopublications and PROV-O have lower detection cost for green-zone content attacks ($O(1)$ per nanopub, $O(|\text{activities}|)$ for PROV-O) but lack native lifecycle governance, so their yellow zone for A5--A7 is larger. Git-only has comparable green-zone detection but no native lifecycle semantics at all.

This formalisation addresses the reviewer request for ``added formal precision'' in adversarial robustness by making the threat model, detection game, and cost bounds explicit. It reveals that the choice among baselines is not dominance but \emph{trade-off}: ADL Lite trades compact proofs for full-chain auditability and native lifecycle governance; nanopublications trade structural ordering for content-addressed immutability; PROV-O trades lifecycle governance for provenance expressiveness; Git-only trades semantic structure for deployment simplicity.

We discuss three classes of threats to the validity of our claims.

\textbf{Internal validity.} The experiments measure architectural correctness (E1--E4, E6, E13--E16) and formal property preservation (E24), not domain-level effectiveness. E2-exhaustive random sampling covers all chains of length up to 3; it does not prove correctness for arbitrary-length chains, although the 100\% pass rate on 10,000 randomized traces (E24, length 2--100) increases confidence. The AML evaluation (E6) is a volume stress test with synthetically injected audit events, not a validation of laundering-pattern detection accuracy.

\textbf{External validity.} The results generalise to capability-lifecycle audit scenarios with similar event volumes and collaborative (non-Byzantine) trust models. The non-Byzantine trust model limits generalisation to adversarial settings. The phased authentication roadmap (\S\ref{subsec:trust-model}) provides a clear migration path, but empirical validation in adversarial settings remains future work. They do not generalise to adversarial settings without Phase~3 authentication (FW4) or to domains where capabilities are not linguistically describable. The evaluation is limited to the AML domain; generalisation to SKILL.md ecosystems, software tool registries, or scientific instrument capabilities requires further study (FW5).

\textbf{Construct validity.} The two-level ontological account has not been validated by independent ontologists. Eight theorems (T1--T7, T9) were machine-verified in Coq 8.18.0; T8 is a Python-level property not in the Coq scope. The randomized trace checker (E24) provided empirical validation. The OWL~2 DL alignment fragment has been validated with ROBOT~\cite{jackson2019robot}: OWL~2 DL profile conformance, structural consistency, and SPARQL constraint verification are confirmed (Section~\ref{app:owl-dl} and Appendix~\ref{subsec:validation-robot}; full report in \texttt{docs/OWL2\_DL\_ROBOT\_VALIDATION\_REPORT.md}). The full operational encoding of $\delta$ and $\gamma$ exceeds OWL~2 DL expressivity and remains future work (FW1). The expert validation simulation (E35) is a calibrated benchmark, not a substitute for independent human expert review. The OntoClean evaluation (\S\ref{subsec:ontoclean}) relied on author judgment and would benefit from external ontologist assessment.

\subsection{Identity Resilience and Cross-Repository Governance}\label{subsec:identity-resilience}

The dual-identifier design (\S\ref{subsubsec:dual-identifier}) addresses a practical problem in multi-agent, multi-repository deployments: how to maintain conceptual identity when EventChains are copied, forked, or archived across organisational boundaries. Three resilience requirements are met. \textbf{R1 (provenance-preserving rehydration):} when a chain is exported from $R$ and imported into $R'$, the genesis hash must remain valid; ADL Lite pins the canonical serialization policy with \texttt{canon\_version} (\S\ref{subsec:crypto}), and the import layer normalises non-canonical formatting (CRLF, float precision) before hash verification---failure to normalise produces a detectable integrity break, not a silent identity change. \textbf{R2 (fork-identity lineage):} a forked child receives a new genesis hash, with the parent-child link recorded as an immutable \texttt{fork-of} relation in the parent chain, preserving lineage without conflating parent and child identities; consumers can trace lineage via the parent's L3 relations or treat the child independently. \textbf{R3 (content-addressed consolidation):} when two agents register semantically identical concepts, the content hash $h_{\text{content}}$ enables near-duplicate detection; the recommended workflow---detect via $h_{\text{content}}$/embedding similarity, emit a \textsc{Relate} \texttt{isomorphic-to}, deprecate the weaker concept---is a manual governance decision (semantic equivalence requires judgment), with $h_{\text{content}}$ providing the \emph{signal} and the \textsc{Relate}/\textsc{Deprecate} events the \emph{audit trail}.

\textbf{Comparison with DID-based identity.} DIDs provide strong actor authentication but weak content binding (they identify \emph{who} made a claim, not \emph{what} it is about). ADL Lite's genesis hash binds identity to the specific event (who + when + what) while the content hash binds to the semantic description; this two-level binding is more resilient than DID-only identity for capability-lifecycle audit because it detects content drift---if a description evolves, the content hash changes, signalling possible re-validation even though the genesis hash (and owning DID) is unchanged.

\subsection{Comparison with Foundational Ontologies}\label{subsec:comparison-foundational}

The ontological analysis in Section~\ref{sec:ontological-analysis} positions ADL Lite as an intermediate between lightweight event-logging systems and heavyweight foundational ontologies. Supplementary Table S27 compares ADL Lite with BFO, DOLCE, and UFO across five dimensions.

ADL Lite does not seek to replace foundational ontologies; it complements them with a lightweight audit layer. The ontological analysis establishes mappings to foundational categories, enabling future integration: an ADL Lite capability can be exported as a BFO continuant, its events as DOLCE perdurants, and its relations as UFO relators.

\subsection{Emerging Standards Integration}\label{subsec:complementary-systems}

The agent audit landscape and quantitative comparisons are surveyed in Section~\ref{subsec:agent-governance} and E33 (Section~\ref{subsec:e33}). Here we identify integration points with emerging standards.

The \emph{Model Context Protocol} (MCP, Anthropic)~\cite{mcp2024anthropic} is implemented as \texttt{mcp\_server.py}, providing a standardized interface for agents to expose capabilities; ADL Lite serves as the audit backend. The \emph{Agent-to-Agent} (A2A, Google)~\cite{a2a2025google} protocol focuses on inter-agent communication; ADL Lite EventChains embed A2A message traces as \texttt{EVIDENCE} events. The \emph{AGNTCY} (Linux Foundation)~\cite{agntcy2025linux} and \emph{NANDA} standards define agent identity and capability ontologies; ADL Lite's L3 predicates align with their relation vocabularies. Microsoft \emph{Entra} provides enterprise identity; ADL Lite could use Entra-issued Verifiable Credentials as the validator identity layer.

\subsection{PROV-O Interoperability Mapping and Loss Analysis}\label{subsec:prov-mapping}

ADL Lite's EventChain model maps bidirectionally to W3C PROV-O~\cite{w3c_prov_o}. Supplementary Table S28 presents the complete bidirectional mapping, and Supplementary Table S29 quantifies the semantic preservation and loss for each mapping direction.

\textbf{Lifecycle event mapping.} Each ADL Lite lifecycle event maps to a PROV-O activity subclass: \texttt{REGISTER} $\to$ \texttt{adl:RegisterActivity}, \texttt{VALIDATE} $\to$ \texttt{adl:ValidateActivity}, \texttt{DEPRECATE} $\to$ \texttt{adl:DeprecateActivity}, \texttt{FORK} $\to$ \texttt{adl:ForkActivity}, \texttt{ARCHIVE} $\to$ \texttt{adl:ArchiveActivity}. Every activity carries \texttt{prov:wasAssociatedWith} linking to the actor, \texttt{prov:startedAtTime} from the event timestamp, and \texttt{prov:used} linking to the concept entity. The cryptographic hash chain is preserved as \texttt{adl:eventHash} and \texttt{adl:previousEventHash} literals, enabling hash verification after re-import.

\textbf{PROV-O profile definition (Turtle).} The normative ADL Lite PROV-O profile---activity subclasses (\texttt{adl:RegisterActivity} through \texttt{adl:ArchiveActivity}) and the datatype properties \texttt{adl:eventHash}, \texttt{adl:previousEventHash}, and \texttt{adl:conceptId}---is provided as a Turtle listing in the supplementary document (Appendix~J, \texttt{supplementary/appendix\_j\_loss\_analysis.tex}).

\textbf{Loss analysis: quantitative preservation metrics.} Supplementary Table S29 quantifies the semantic coverage of the bidirectional mapping. Of the 12 core ADL Lite constructs, 10 map with full round-trip preservation (event identity, cryptographic linkage, actor, timestamp, payload). Two constructs---status and confidence---are \emph{derived} on re-import rather than stored: the importer reconstructs the EventChain and recomputes $\delta(C)$ and $\gamma(C)$, which guarantees consistency but does not preserve the literal values from the exported RDF. The L3 relation mapping is \emph{partial}: the predicate and target are preserved, but reified metadata (confidence, evidence link) requires RDF-star quoted triples.

\textbf{Unmapped semantics: scope differences.} Five ADL Lite features have no direct PROV-O equivalent, reflecting deliberate scope differences rather than implementation gaps: (i)~\emph{deterministic state derivation} ($\delta$, $\gamma$ are application-layer semantics); (ii)~\emph{precondition rules} (PROV-O has no state-transition constraints); (iii)~\emph{fork semantics} (\texttt{prov:wasDerivedFrom} models lineage but not fork-determinism or child-chain independence); (iv)~\emph{confidence aggregation} (PROV-O records values but has no aggregation mechanism); and (v)~\emph{graded weakening} (PROV-O's \texttt{prov:wasInvalidatedBy} is binary). These are documented as extension points for future PROV-O profile definitions.

\textbf{Round-trip validation and import path.} The bidirectional mapping is implemented in \texttt{prov\_export.py} / \texttt{prov\_import.py} and evaluated by E12: every exported EventChain parses as valid Turtle, re-imported chains have identical \texttt{event\_id}s and hashes, recomputed status/confidence match the original derived values, and L3 relations are preserved under RDF-star export (all four properties hold on 201 AML EventChains, 9,300 events, zero failures). External PROV-O traces are imported by grouping \texttt{prov:Activity} instances by \texttt{adl:conceptId}, sorting by \texttt{prov:startedAtTime}, converting each activity to an Event, recomputing hashes, and running \texttt{verify\_integrity()}---enabling ADL Lite to consume provenance from existing PROV-O pipelines (e.g., workflow systems) and audit them through $\delta(C)$ and precondition rules.

\subsection{Positioning and Distinctive Properties}\label{subsec:positioning}

ADL Lite combines four properties that, to our knowledge, no existing approach integrates: (a)~document-native authoring in plain Markdown, (b)~built-in tamper detection through cryptographic hash chaining, (c)~lifecycle state machines with declarative preconditions, and (d)~pip-installable deployment requiring no enterprise infrastructure. This enables multiple LLM agents to propose, challenge, and refine concepts within a shared document-native ontology, with the EventChain providing the audit trail and deterministic derivation mechanism that renders decentralized authorship accountable.
